% chapters/ch07.tex —— 第 7 章
% 体例：具体失败 → 设计问题 → 机制解释 → 对照实验 → 条件变化 → 选择原则
% 数字身份纪律：本章外部数字全部来自方法卡与原文抽取所登记的节号/表号，一律 \cited{}；
% 由引用数字算出的量用 \calcnum{}；本书实测（E02 的自评与环境结果对照）用 \measured{}。
% 本章无本轮新做的实测。E04 属 B 级受限实验，尚未运行，本章不报告它的任何数字。
\chapter{指令与上下文的自动优化}
\label{ch07}

\chaptergoal{怎样把一次失败变成一次有效的策略修改？}
\chapterprops{优化收益怎样推广到未参与选择的任务？本章分别检查搜索对象、诊断依据与预算分配，区分内容丢失、选择偏差与泛化不足；E04保留为后续受限研究设计。}
\chapterfig{图~\ref{fig:ch07-context}：上下文更新的两条路径与内容检查位置}
\chapterpitfall{把提示变短当成收敛；把上下文变长当成信息在增长；
把优化集上的提升当成能力变化}

% ---------- 1. 具体失败 ----------
\begin{failurebox}
在一个交互式 Agent 基准上，有人让语言模型在每一步把积累下来的上下文整体重写一遍。
到第 \cited{60} 步时，这份上下文有 \cited{18{,}282} 个 token，准确率 \cited{66.7\%}；
到第 \cited{61} 步，它被重写成 \cited{122} 个 token，准确率掉到 \cited{57.1\%}——
\emph{低于完全不做适配时的} \cited{63.7\%}\cite{ace-2510}。

一步之间丢掉 \calcnum{99.3\%} 的上下文，不是模型突然变笨，
这是整体重写可能丢失既有内容的一个案例，不能据此断言所有重写必然失败：
模型的输入是整份上下文，它输出的是它认为重要的东西，
若系统没有保存前后版本及差异，被省略的内容就难以追查。

ACE 还讨论了简短泛化指令可能遗漏任务细节的风险；其中测试生成案例转述自另一项研究，本章不将转引视为本书独立核验的实验。后文将用 GEPA 原论文实际展示的第二跳检索提示，分析有价值的修改具体改变了什么。
\end{failurebox}

\designquestion{判据不是``上下文有没有变化''，而是：
\textbf{被删掉的那部分，有没有一个单独的动作对它负责}。
明确的删除事件便于审计；整体重写也可以通过前后版本、条目映射和差异检查获得可追溯性。}

第 6 章处理``记忆操作由什么决定''，本章处理它的近邻：
当被优化的对象从``一条记忆''换成``一整份上下文或一段指令''时，
同样四个动作（写、留、取、删）里，``留''和``删''的账会变得难算得多——
因为一次重写可以同时完成全部四个动作，而账上只显示一行``上下文更新了''。
本章要把这一行拆回四个动作。

% ---------- 2. 机制解释 ----------
\section{搜索空间：指令、示例与上下文选择}
\label{sec:ch07-space}

把``优化''这个词拆开，第一个要回答的问题是\emph{优化的对象是什么}。
三条主线给出了不同答案，搜索规模取决于候选数、模块数和组合方式。

\textbf{指令与字段描述。}DSPy 把一次语言模型调用声明成 signature（用自然语言声明输入、输出字段及其语义），
module 把 signature 与若干 demonstrations 封装在一起，
teleprompter 是那个负责改它们的优化器\cite{dspy-2310}。
原论文将典型优化器概括为三个阶段，并非每种优化器都必须依次执行：先为每个模块生成指令、字段描述与演示的候选，
再用随机搜索或 TPE 在离散候选上选优，最后做 ensemble 与控制流这类高阶程序优化\cite{dspy-2310}。
搜索对象包括模块指令、字段描述与示例等可编译内容。参数槽数可以随模块数增加，但组合候选数可能按各模块候选数的乘积增长，不能笼统称为线性。

\textbf{示例。}DSPy 原论文重点研究 demonstrations：
用高温采样跑出多阶段轨迹，只保留整条轨迹通过指标的那一些\cite{dspy-2310}。
这一步是拒绝采样，它的筛选单位是一条完整的轨迹，
而不是``这条轨迹里的哪一步起了作用''。

\textbf{上下文选择。}ACE 把上下文做成一组条目（bullet），
每条由元数据（唯一标识、被标记为有用/有害的计数）与内容两部分构成\cite{ace-2510}。
因此可以分别讨论条目的生成、修改、合并及读取；它不只是一个固定条目集合上的检索选择问题。
ACE 报告的三个基准上的整体结果见表~\ref{tab:ch07-ace}。

\begin{table}[htbp]
\centering
\caption{ACE 原论文中按适配条件分开的结果摘录}
\label{tab:ch07-ace}
\begin{tabular}{>{\raggedright\arraybackslash}p{3.3cm}>{\raggedright\arraybackslash}p{3.4cm}ccc}
\toprule
基准与条件 & 方法 & 原文平均 & FiNER & Formula \\
\midrule
AppWorld 基础 & ReAct & \cited{42.4} & --- & --- \\
AppWorld 离线有标签 & ICL / GEPA & \cited{46.0 / 46.4} & --- & --- \\
AppWorld 离线有标签 & ACE & \cited{59.4} & --- & --- \\
AppWorld 在线无标签 & DC / ACE & \cited{51.9 / 59.5} & --- & --- \\
金融基础 & Base LLM & \cited{69.1} & \cited{70.7} & \cited{67.5} \\
金融离线有标签 & GEPA & \cited{72.5} & \cited{73.5} & \cited{71.5} \\
金融离线有标签 & ACE & \cited{81.9} & \cited{78.3} & \cited{85.5} \\
金融在线有标签 & ACE & \cited{76.6} & \cited{76.7} & \cited{76.5} \\
金融在线无标签 & ACE & \cited{72.9} & \cited{67.3} & \cited{78.5} \\
\bottomrule
\end{tabular}

\vspace{2mm}
\begin{footnotesize}
数字为原论文 Table 1、Table 2 的百分数读数\cite{ace-2510}。AppWorld 平均汇总 TGC 与 SGC 在两个划分上的值，金融列为准确率，不能混称统一成功率。在线 AppWorld 的 ACE 含离线 warmup，见原文 Table 3。这里只在相同条件内解释比较，不把各行串成方法排名。
\end{footnotesize}
\end{table}

\subsection{同一框架在两个任务上，反馈作用可以不同}

ACE 的金融结果提供了一个比“没有标签也能适应”更具体的观察。在线无标签条件下，FiNER 从基础模型的\cited{70.7}变为\cited{67.3}，Formula 则从\cited{67.5}变为\cited{78.5}\cite{ace-2510}。按原表读数相减，分别是\calcnum{-3.4}和\calcnum{+11.0}个百分点；两任务平均改善\calcnum{3.8}个百分点，并不代表每一个任务都受益。这些差值只是原文点估计的算术转换，本书没有其逐题数据，不能据此补造显著性结论。

无标签不等于无信息：执行返回值、工具状态和环境变化仍可能提供约束。原文将任务差异与可靠反馈的可用性联系起来，表格却没有单独操纵所有因素，因而不能据此估计反馈可靠性的纯因果效应。

\section{诊断与解释：谁有资格说``错在哪''}
\label{sec:ch07-diagnose}

第二个问题是：那句``下次应该这样做''是谁写的。

ACE 把更新流程分成三个角色：Generator 产生推理轨迹，Reflector 从成功与失败里抽具体的洞见，
Curator 把这些洞见合成结构化的增量条目\cite{ace-2510}。
分出来的关键是\emph{评估与洞见抽取}和\emph{策展}不是同一个模型调用：
Reflector 提炼成功与失败的具体经验，Curator 提出条目更新；
最终增量由确定性的合并程序应用。不能把反思者的建议直接当成已经完成的发布\cite{ace-2510}。

第 4 章研究了与此相关的另一个问题：反馈从哪里来。
本书 E02 在任务集 v1 上比较了四个主要反馈配置，另有噪声对照与无记忆基线。其中``内在自评''（同一后端在独立评审请求中按自评措辞判断代码）
的封存题回合通过率是 \measured{0.6125}，``环境结果''（删去验收总结行的执行输出）是 \measured{0.7000}；后者不显式添加总判定，但回溯和断言细节仍可能透露成败。
两臂通过率之差 \measured{-0.0875}，95\% 区间 \measured{[-0.2625,+0.0875]} 含 0。两个反馈臂分别训练并构建经验，比较对象是完整反馈方案，不能把差值单独归给某句反馈措辞或标签强度。
\begin{itemize}
\sloppy
  \item 这句话的边界要说清楚：E02 的自评臂做的是\emph{给出一个通过与否的判定}，
        反馈判定随后还会进入经验构建；反馈判断与最终经验文本仍是不同环节。
  \item E02 的评审与自评请求共享同一后端，载荷包含代码与测试文件，不继承求解历史；差异主要在提示措辞。该比较尚不足以归因到“自己”这个角色。
  \item 两个实验的任务内循环不同。E02 归档提示未列出验收工具，轨迹中该工具调用为零，测试文件可读；保留代码中的分派入口尚未与历史运行作源码绑定。两次实验的绝对数字不可直接相减。
\end{itemize}

ACE 还报告了不使用标注监督、借助自然执行反馈的适配\cite{ace-2510}。
它与第 4 章的层级表共同提示：
Reflector 可以将执行信息整理为具体建议，但不能凭解释的细致程度提高原始反馈的真值可靠性。翻译这一步做没做，是可以检查的——
看产出的条目里有没有指出位置或条件。

\subsection{沿着一次失败修改：第二跳检索的例子}
\label{sec:ch07-twohop}

GEPA 原论文 Figure 2 展示了一次具体的提示变化：被修改的是多跳问答系统的第二跳查询模块，输入为原始问题和第一跳摘要，输出为新的检索查询\cite{gepa-2507}。初始指令只要求根据输入产生查询；优化后的指令则解释了第二跳的职责，要求寻找第一跳尚未覆盖、但回答问题所必需的信息，并提醒模型不要只是重述原问题。这个例子适合用来理解提示优化，因为变化能够定位到一个模块的输入输出，而不只是“请更认真”一类笼统要求。

这里应先区分原文展示与本书解释。论文给出了初始提示和优化提示，但单凭这张图，不能重建每一轮候选修改的全部搜索历史，也不能证明最终提示中的每一句都对测试收益有独立贡献。下面沿信息流分析它为什么可能有用，是机制解释；若要估计每一条规则的边际作用，还需单独消融。

设原问题需要先找出某位作者的出生地，再确定该地点所属的行政区域。第一跳已经返回出生地，如果第二跳仅把原问题换一种说法再搜，可能再次得到同一份人物介绍。这里的失败不是问题太短，也不一定是检索器太弱，而是第二跳没有利用第一跳结果缩小剩余的信息缺口。一个有针对性的修改应让模块先识别“已有事实”和“尚缺事实”，再构造下一条查询。

这是一段\textbf{本书构造的任务解释}，不是论文中额外报告的题目或运行结果。它对应的教学提示可以写成：先从摘要中列出与问题有关且已有依据的实体关系；再识别回答仍缺的一项关系；只为缺失关系生成查询，保留必要的实体限定；若摘要没有可靠实体，不得凭空补造，应返回需要重新获取的输入。读者可以看到，修改既增加了正向动作，也增加了适用前提和回退行为。

“不要重复第一跳”本身还不够。第一跳可能检索错了对象，也可能返回空结果；此时绝对禁止重复查询会阻止合理的修复。优化提示需要说明，避免的是在已有充分信息时无意义重复，不是禁止对失败查询进行改写和重试。这种限定使提示从一句口号变成可检验的行为要求：当第一跳信息不足时允许恢复，当信息充足时应转向缺口。

\subsubsection{诊断要对应可改动的责任范围}

假如最终答案错误，第二跳提示不一定是合适的修改位置。错误可能发生在第一跳摘要，将甲的出生地写成乙的；也可能发生在第二跳检索，搜索词正确但返回资料过时；还可能发生在最终回答阶段，正确证据已经存在却被错误组合。把所有失败都交给第二跳模块，会逼它用更长的提示补偿自己无法控制的上游问题。

因此，诊断材料应沿模块边界展开：第二跳实际收到什么、生成什么、工具返回什么、最终回答用了哪些证据。若输入摘要已经包含错误实体，候选修改应针对摘要或输入校验；若查询遗漏唯一限定条件，才适合改查询提示；若生成结果违反输出格式，则可以先修格式约束。自然语言反思能帮助定位这些问题，但它提出的归因仍需与轨迹对照。

在 DSPy 的抽象下，signature 给出输入输出语义，module 实现相应推理方式，优化器选择指令或示例等参数\cite{dspy-2310}。这种结构使错误定位有了可操作的边界，却没有自动解决语义正确性。字段名是 \texttt{summary} 不保证它忠实，输出能被解析为查询也不保证它找到了缺失事实。格式验收与任务验收应当分别保留。

一个候选修改可以附带简短的“诊断—干预—预期变化”记录。诊断是第一跳内容被重复检索；干预是要求第二跳显式识别未覆盖关系；预期变化是查询转向该关系，而非简单重述问题。评价时先检查预期的行为是否发生，再看任务表现。如果查询确实改变但最终得分未改善，说明干预执行成功而效果未获支持；如果查询没有改变，则不能直接说机制无效，可能是提示没有被执行或被其他内容覆盖。

\subsubsection{正确示例也可能携带错误理由}

DSPy 的自举可以从通过最终指标的多阶段轨迹中收集示例\cite{dspy-2310}。这种做法减少了逐步人工标注需求，但终点通过并不保证中间每一步都适合复用。一个回答可能靠后续检索纠正了错误摘要，也可能偶然命中答案。如果把错误摘要当成标准的模块输出，下一轮优化就可能把“最终侥幸成功”变成局部行为规范。

对多跳问答而言，需要检查的不是所有推理句子是否漂亮，而是被选为示例的中间产物是否符合其接口契约。摘要应包含可由输入支持的事实，查询应指向未解决的信息，引用应对应实际证据。检查可以是确定性的，也可以由独立评价完成；若只能验收最终答案，就应承认中间示例的质量没有被单独确认。

示例顺序也属于候选的一部分。同一组演示改变排列，可能改变模型采用的模式；加入一个很长的示例，还可能挤掉最后的限制条件。只保存示例集合而不保存最终渲染文本，不足以复现提示程序。对于多模块系统，应保存每个模块实际使用的顺序、字段描述、模型标识和解析配置，而不只保存顶层说明。

\section{提示程序的契约：可搜索部分与固定部分}
\label{sec:ch07-contract}

提示优化不应默认允许修改所有输入。任务的输出格式、工具权限、预算和证据要求，可能是产品或实验的固定条件；策略说明、演示和检索提示，才是本轮允许搜索的部分。若优化器能够把“必须提供证据”改成“不需要证据”，它可能提高某个宽松指标，却已经改变了任务。

一个实用的划分是把最终请求分成固定契约、候选内容和任务数据。固定契约由框架加载，候选内容有版本标识，任务数据按明确边界插入。评价时保存三者渲染后的实际请求，因为相同片段在不同位置和角色下可能产生不同效果。这个划分并不提供理论上的绝对隔离，但使“优化了什么”成为可核对的问题。

固定部分也应尽量小。把所有设计偏好都锁进不可改动区域，会使优化器只能在无关措辞上搜索；完全不设边界，则可能让它通过改变验收标准获得高分。应从研究问题出发决定哪些改变算合法干预。例如，若比较候选选择算法，应固定输出解析器；若研究解析器与提示的共同设计，则可以允许它变化，但要将最终程序作为完整产物评价。

提示文本中的指令与引用材料还需要区分。错误日志、网页内容、用户提供的文档或旧经验里可能出现命令式句子，它们首先是待处理的数据，不能因为进入上下文就获得修改系统规则的权力。优化器从失败轨迹提炼经验时，应保留这层来源身份。否则一次外部材料中的伪指令，就可能被概括成下轮长期生效的策略。

这种风险不只发生在明显恶意输入中。测试文件里的建议、工具错误消息中的示例命令、旧版本的运行说明，都可能被误认作当前任务要求。第5章的经验审计已经出现材料与实际动作混淆的迹象。提示优化同样需要区分“轨迹里出现过一句话”与“这句话是合法的行为约束”。生成的候选应能够追溯到具体观察，而不只是引用整条轨迹作为模糊依据。

\subsection{可编辑范围决定搜索规模}

假定一个程序有若干模块，每个模块允许从几个指令候选中选择。即使模块参数的数量随模块数线性增长，全部组合的数量也可能按候选数相乘。再加入示例子集、示例顺序和控制流，搜索空间会进一步扩大。模块化提高了可描述性，不意味着优化问题天然变小。

这解释了为什么不能把所有改动一次性放开。在预算有限时，可以先只允许一个模块变化，用其余固定模块判断候选是否产生预期行为；之后再开放相邻模块，检查组合是否仍有效。这是缩小解释范围的一种工程做法，不是保证获得全局最优的方法。逐模块优化可能错过只有联合改变才有效的组合，因此需要在结果中说明这一搜索限制。

控制流变化尤其要单列。新增一个验证模块，可以提高成功率，但也增加了调用次数；在错误时重试，可以改善任务结果，却改变了执行策略。若这些改变被隐藏在“提示更好了”之下，读者就无法分辨收益来自语义指导、更多计算，还是更好的失败恢复。DSPy 支持更高阶程序优化，本书对提示优化的评价必须保留这一边界，不能将所有系统级收益归给一段字符串。

\subsection{候选差异应能转成一组可检查的问题}

保存文本差异只是第一步。一次修改可能把“只在证据不足时再次检索”改成“再次检索”，字符删除很少，行为范围却发生了显著变化。另一种修改只是替换同义表达，字符差异很大，但预期行为几乎不变。若用编辑距离衡量修改幅度，应明确它只是文本变化量，不是策略变化的通用尺度。

对每个重要候选，可以围绕条件、动作、禁止项和回退行为建立检查问题。条件是否从特定任务扩成全部任务；动作是否增加了工具调用；禁止项是否与正常恢复冲突；回退是否会跳过必要的验收。这些问题能把自然语言修改变成可审读对象，也有助于构造少量有针对性的回归任务。

回归任务不能全部来自触发修改的同一道失败题。否则候选只需补上一句该题的答案，就能看似修复。更有判别力的是保留同一机制、改变表面信息的任务，再加入一个不应触发该规则的反例。它们可以用于开发期筛选；最终泛化报告仍需独立任务，不能把针对候选编写的检查重新包装为未见测试集。


\section{候选生成、选择与预算分配}
\label{sec:ch07-budget}

第三个问题是钱花在哪。这里有一个容易被误读的对照。

GEPA 采样轨迹、用自然语言反思诊断问题、提出并测试提示更新，
再从自己尝试过的 Pareto 前沿上组合互补的经验\cite{gepa-2507}。
按原论文 Table~1 的 Qwen3 8B 结果，GEPA 与 GRPO 的六任务汇总分分别为 \cited{54.85} 与 \cited{48.91}，相差 \calcnum{5.94} 个分数点；HotpotQA 一列相差 \calcnum{19.00} 个分数点。它并非每项都领先：AIME-2025 上为 \cited{32.00} 对 \cited{38.00}\cite{gepa-2507}。同表 GEPA 相对 MIPROv2 的汇总差是 \calcnum{7.01} 个分数点，AIME-2025 差是 \calcnum{12.00} 个百分点。不能把摘要中“超过10\%”直接改写成汇总提升超过十个百分点，也不能把异质任务汇总分一律当成准确率。

原文所谓最多减少约 \cited{35} 倍 rollout，还需要区分找到最佳候选的时点与全部搜索投入。Table~1 的说明以 IFBench 在 \cited{678} 次 rollout 时找到最佳提示，对照 GRPO 的 \cited{24000} 次；但同表标作 GEPA (+Merge) 的预算行给出的 IFBench 总优化预算为 \cited{3593} 次\cite{gepa-2507}。前者可描述观察到的最佳候选出现时点，不能替代整个运行的最终账单。若事后才知道哪一步最好，也不能假设部署者事先知道可以在那一步停止。

\begin{principlebox}
\textbf{这条对照的正确读法。}
它不是``提示优化优于强化学习''。
它报告的是特定任务、模型与优化配置下的样本效率差异；这不单独隔离自然语言反馈相对于标量反馈的因果贡献。
这里比较的资源单位是 rollout，不能直接换算为总 token、墙钟或美元。
约 \cited{1/35} 比较的是最佳提示出现时点的 rollout 数与 GRPO 预算，并非两次完整优化的总预算比。
两边单次 rollout 的资源消耗也不保证相等，因而不能读成``便宜 35 倍''。
\end{principlebox}

代价这一侧 ACE 按场景报告了适配成本（表~\ref{tab:ch07-cost}）。

\begin{table}[htbp]
\centering
\caption{适配成本对照：离线（AppWorld，对 GEPA）与在线（FiNER，对 DC）}
\label{tab:ch07-cost}
\begin{tabular}{>{\raggedright\arraybackslash}p{2.6cm}cc>{\raggedright\arraybackslash}p{4.0cm}}
\toprule
场景 & 方法 & 延迟（秒） & rollout数 / token费用 \\
\midrule
离线 & ReAct + GEPA & \cited{53898} & \cited{1434} 次 \\
离线 & ReAct + ACE & \cited{9517} & \cited{357} 次（\cited{-75.1\%}） \\
在线 & DC (CU) & \cited{65104} & \cited{17.7} \mbox{美元} \\
在线 & ACE & \cited{5503}（\cited{-91.5\%}） & \cited{2.9} \mbox{美元}（\cited{-83.6\%}） \\
\bottomrule
\end{tabular}

\vspace{2mm}
\begin{footnotesize}
数字身份：[引用] 原论文 Table~4(a)(b)\cite{ace-2510}。
这些金额为论文在其服务与计价条件下的历史报告，不能作为本书实验费用或当前服务报价。
\end{footnotesize}
\end{table}

原论文附录 A.3 另有一组细粒度适配统计，明确将 ACE 设为一次 epoch、一次反思改进，GEPA 使用官方实现的 heavy 配置。该组输入 token 与输出 token 相对 GEPA 分别降 \cited{80.8\%} 与 \cited{83.6\%}，
原文归因包括减少提示验证循环及用局部更新替代整体重写\cite{ace-2510}。
该比较没有将这些因素逐一隔离，不能仅凭总下降比例判定哪一项贡献最大。
更要注意，该组 Table~12 的总 rollout 是 ACE \cited{2075} 对 GEPA \cited{1455}，增加 \cited{42.6\%}，与表~\ref{tab:ch07-cost} 的 rollout 下降不是同一组统计。正文表4与附录的汇总不能互换，附录数字也不能解释为较少调用带来的同一笔节省。它说明调用次数、单次请求长度与组件组成必须一起看\cite{ace-2510}。

若要区分更新方式与候选评估预算的作用，可以尝试在固定模型、任务、可见反馈和成本范围下交叉安排两项因素，而不是只比较两个完整方法的端点。组合实现若需要改动原系统，应先检查每个单元确实实现了预定因素，并报告额外适配。两项因素还可能交互：局部更新只有配合特定评估流程才产生节省时，总下降比例就不能唯一拆成两份贡献。这是本书提出的对照设计，不是上述原表已经完成的消融；第\ref{ch15}章的四格交互给出同类比较的计算方式。

\subsection{预算相等，要相等到哪一层}
\label{sec:ch07-budget-scope}

一次 rollout 可以包含多次模型请求、检索和工具调用。GEPA 的问题定义把系统运行及评价作为受预算约束的单位\cite{gepa-2507}；这有利于讨论如何从少量任务试验中获得改进，但它与服务账单上的调用次数不是同一个量。多模块程序、长反思文本和短分类任务，可能拥有相同 rollout 数而消耗完全不同的资源。

比较提示优化方法时，至少要说明生成候选的模型、评估候选的模型和最终执行模型是否相同。使用更强的反思模型指导较弱执行模型，可能是合理的系统设计，但教师的知识和计算不能隐形。若只按执行模型的调用数宣称高样本效率，应同时披露教师调用及其可见的反馈范围。

等预算可以有不同定义。固定评估任务数，关注如何利用同样多的任务反馈；固定总 token，关注语言模型资源；固定费用，关注具体计价条件下的支出；固定延迟，关注用户等待。它们往往无法同时完全相等，因此不应只写“公平比较”而不说明约束。可以选择一个主要预算，再完整报告其他资源结果。

查询预算还需要包括失败与拒绝。解析失败后再请求一次候选，是搜索成本；评估超时后的重跑，也是成本；模型服务重试可能增加请求尝试数而不产生新的有效评价。应用层、网络层与实验层的事件应分开记录，避免将没有成功返回的运行从分母和成本中一起删除。第8章的技能成本账会再次遇到这个问题。

提示长度有两种身份。最大可用长度可以是预设预算，实际生成长度则是方法结果。若研究“在相同长度上限下谁更有效”，就应固定上限；若研究“多长的上下文值得支付”，则应改变长度并报告质量与费用关系。两种问题都成立，关键是不要在看到结果之后，才选择最有利的比较口径。

\subsubsection{缓存改变账单，不改变信息量的定义}

ACE 在其 GPT-5.1 缓存研究中报告了较高的输入缓存比例和相应费用下降\cite{ace-2510}。这说明长上下文的原始 token 数不一定按未缓存价格逐次收费，却不表示这些 token 没有进入模型，也不表示它们不会影响注意与任务行为。成本与信息干扰仍需分别评价。

缓存收益还取决于请求排列与可复用前缀。若每次把新增条目插到最前方，稳定前缀可能缩短；若将固定规则放前面、变化材料放后面，可能更有利于复用，但具体命中仍要由服务记录确认。并发、缓存生命周期与后端策略也会影响结果。书中引用论文的历史缓存测量时，应保留模型与实验条件，不把它推广为任意部署的保证。

还应避免把费用降低写成 token 消失。未缓存 token、缓存 token 和输出 token 可能使用不同单价，实际费用是这些数量按相应价格加总的结果。若论文只报告费用比例，不能逆推出上下文实际缩短了同样比例。资源表最好同时保留数量、单位价格及适用日期；本章不以历史论文报价替代当前服务价格。

\subsubsection{优化成本如何分摊到部署}

一次离线优化可以产生长期复用的提示，因此需要区分开发投入和部署费用。若自动搜索花了较多资源，却在大量后续任务中稳定减少调用，它可能在约定使用量下值得投入；如果提示只用一次，开发投入就难以摊薄。这里与技能固化相似，但复用资产是提示程序，不是可执行代码。

分摊时必须固定资产版本及复用范围。提示每用几次就再次优化，后续搜索不应被藏进第一轮的已付成本；换模型后重新编译，也是一轮新的适配投入。若报告的是预期部署量，可以提供情景分析；若报告已实现费用，应只按实际运行次数计算。把预期百万次使用写成当前已经获得的摊销收益，会夸大方法的经济性。

从这一角度看，固定提示基线并不是“没有投入”。它可能来自已有人工设计或先前优化。实验可以将共同起始提示视为共享前置条件，但应声明其来源；若某一臂独享大量历史试错，就不能同时把它称为零搜索基线。比较范围从哪里开始，决定了哪些费用可以共同排除，哪些必须归到特定方法。


\section{结构化增量更新：膨胀、冲突与信息丢失}
\label{sec:ch07-delta}

% figures/ch07-context.tex —— 第 7 章核心图（全书引用入口）
% 用法：在 book/chapters/ch07.tex 中 % figures/ch07-context.tex —— 第 7 章核心图（全书引用入口）
% 用法：在 book/chapters/ch07.tex 中 % figures/ch07-context.tex —— 第 7 章核心图（全书引用入口）
% 用法：在 book/chapters/ch07.tex 中 \input{figures/ch07-context}
% 图主体在 ch07-context-body.tex，便于单独编译与复用。
\begin{figure}[htbp]
\centering
\input{figures/ch07-context-body}
\caption{上下文更新的两条路径：整体重写与增量提案都需要内容检查、版本记录和外部评价。更新形式改变审计方式，不自动保证正确性。}
\label{fig:ch07-context}
\end{figure}

% 图主体在 ch07-context-body.tex，便于单独编译与复用。
\begin{figure}[htbp]
\centering
\begin{tikzpicture}[
 box/.style={draw,rounded corners=2pt,text width=31mm,minimum height=13mm,align=center,font=\small},
 arrow/.style={-{Stealth[length=1.6mm]},thick}
]
\node[box,fill=gray!10] (old) at (0,0) {原上下文\\保存版本};
\node[box,fill=orange!12] (rewrite) at (4.2,1.15) {整体重写\\生成完整候选};
\node[box,fill=blue!10] (delta) at (4.2,-1.15) {增量提案\\标明目标与修改};
\node[box,fill=orange!8] (diff) at (8.4,1.15) {比较前后差异\\检查遗漏与新增};
\node[box,fill=blue!8] (merge) at (8.4,-1.15) {检查版本与冲突\\合并合法提案};
\draw[arrow] (old)--(rewrite);
\draw[arrow] (old)--(delta);
\draw[arrow] (rewrite)--(diff);
\draw[arrow] (delta)--(merge);
\node[align=left,text width=136mm,font=\footnotesize] at (4.2,-2.5)
{两种路径都应保存候选、检查内容变化并独立评价。增量便于定位修改，仍可能追加错误或造成冲突；整体重写有遗漏风险，也可以通过版本差异审计。图中没有性能排名。};
\end{tikzpicture}

\caption{上下文更新的两条路径：整体重写与增量提案都需要内容检查、版本记录和外部评价。更新形式改变审计方式，不自动保证正确性。}
\label{fig:ch07-context}
\end{figure}

% 图主体在 ch07-context-body.tex，便于单独编译与复用。
\begin{figure}[htbp]
\centering
\begin{tikzpicture}[
 box/.style={draw,rounded corners=2pt,text width=31mm,minimum height=13mm,align=center,font=\small},
 arrow/.style={-{Stealth[length=1.6mm]},thick}
]
\node[box,fill=gray!10] (old) at (0,0) {原上下文\\保存版本};
\node[box,fill=orange!12] (rewrite) at (4.2,1.15) {整体重写\\生成完整候选};
\node[box,fill=blue!10] (delta) at (4.2,-1.15) {增量提案\\标明目标与修改};
\node[box,fill=orange!8] (diff) at (8.4,1.15) {比较前后差异\\检查遗漏与新增};
\node[box,fill=blue!8] (merge) at (8.4,-1.15) {检查版本与冲突\\合并合法提案};
\draw[arrow] (old)--(rewrite);
\draw[arrow] (old)--(delta);
\draw[arrow] (rewrite)--(diff);
\draw[arrow] (delta)--(merge);
\node[align=left,text width=136mm,font=\footnotesize] at (4.2,-2.5)
{两种路径都应保存候选、检查内容变化并独立评价。增量便于定位修改，仍可能追加错误或造成冲突；整体重写有遗漏风险，也可以通过版本差异审计。图中没有性能排名。};
\end{tikzpicture}

\caption{上下文更新的两条路径：整体重写与增量提案都需要内容检查、版本记录和外部评价。更新形式改变审计方式，不自动保证正确性。}
\label{fig:ch07-context}
\end{figure}


面对本章开头那类塌缩风险，ACE 将三项设计配合使用：
独立的 Reflector 提炼经验，增量 delta 用局部更新替代整体重写，
grow-and-refine 则将持续扩张与基于条目标识、计数和语义去重的整理配对\cite{ace-2510}。

三条里第三条最容易被忽略。
只有净追加超过合并和删除时，上下文才会增长；是否设置准入门槛本身不能决定长度的变化。
ACE 将新增条目的追加与既有条目的更新分开，并用语义去重裁剪冗余；原文允许每次增量后主动整理，也允许到达长度限制时再整理。
而对一条已经被判定为有害的条目，ACE 的做法是给它记上计数\cite{ace-2510}：
计数本身不删除任何东西，但它把``这条可能是坏的''变成了一个可查的量。

\begin{principlebox}
\textbf{增量更新换了什么。}
整体重写把四个动作（写、留、取、删）压缩成一次生成，
如果没有保存前后差异，就难以解释单条内容为何消失。
在 ACE 描述的合并方式中，未涉及的条目得以保留，新标识对应追加，既有标识对应原位更新；语义去重另行处理冗余裁剪。
工程上还可以设计点名移除操作并单独记账，但不能仅凭 delta 这个名称，就推断实现已经支持删除指令或验证了删除安全性。实际审计仍依赖日志、版本与合并规则，延迟和 token 收益则要另行测量。
\end{principlebox}

代价也在同一处。若追加超过去重和淘汰，上下文会增长，读取与维护负担也可能上升；
这与第 8 章技能库增长的账是同一笔账，只是对象从代码换成了文本。
ACE 的处置是 grow-and-refine 里的 refine 那一半，
有用/有害是同一条目的标记，不直接表示两条条目之间的矛盾关系。原文附录 A.4 还说明整理器可按有害标记过滤条目，A.6 讨论过时、低效内容的裁剪；这些维护不能简化为只追加或只增加计数，也没有使任意局部修改自动一致\cite{ace-2510}。

\subsection{增量更新之后，仍然需要判断改变是否合法}
\label{sec:ch07-update-contract}

将上下文拆成条目以后，每次更新可以只提交新增或修改的部分，减少重复生成整份文本的负担。ACE 的结构化条目与非模型合并逻辑提供了这样的实现方向\cite{ace-2510}。但是，增量形式解决的是更新的定位和执行方式，不能自动保证新增内容正确，也不能保证两个局部正确的修改组合后仍然一致。

设已有条目规定“只有证据不足时再次检索”，新条目根据一次失败提出“回答前务必再次检索”。二者单独都能解释某些案例，合在一起却给出不同的触发条件。简单追加会把冲突留给执行模型；基于相似度删除其中一条，又可能删除了更精确的限定。策展需要判断两条规则之间是重复、补充、特例还是矛盾，而不只是比较文本距离。

这也是为什么有用与有害计数不能独自承担维护。某条规则在常见任务上频繁被标记有用，并不排除它在罕见但重要的场景中有害；某条几乎不被调用的例外，也可能是防止严重错误的必要条件。计数需要关联发生的任务与触发条件，才能帮助解释作用范围。将所有任务的标记直接相加，会让高频任务支配低频任务的记忆选择。

\subsubsection{一次 delta 应有可检查的前置条件}

一个增量提案可以指定目标条目、期望版本、操作类型、修改后的内容和证据来源。新增与修改应分开：若提案引用的条目已被其他更新改变，合并器应报告冲突，不能把它当作仍然存在的旧条目直接覆盖。若提案只增加计数，则应明确计数对应的是一次新的评价事件，避免重放同一事件造成虚假热度。

并行提出的两条 delta 只有在不争用相同条目、不改变相互依赖的条件时，才容易直接合并。两个提案都修改同一条规则，或者一个删除另一个所引用的条目，就需要确定顺序或重新审查。ACE 的局部更新支持批量合并，这不等于任意文本修改都可交换。工程实现应识别依赖，而不能由“并行”一词推导没有冲突。

整体重写也可以获得类似的检查。先保存原版本，将新文本中的规则映射回旧条目，列出保留、修改、移除和新增，再检查被移除的条件是否仍由其他内容覆盖。这个过程可能比显式 delta 更贵，也可能在语义对齐时产生不确定性，但它说明可审计性来自记录和验证流程，而不只来自更新形式。

当条目关系无法可靠判断时，可以让候选保持未应用状态，继续使用旧版本完成任务。暂不应用不等于认定候选有害，而是承认当前证据不足。对于必须立即处理的任务，应由预定的回退规则决定使用哪些已验证内容；不宜让策展器为了完成本轮更新，强行把尚未解决的矛盾改写成统一结论。

\subsubsection{条目内容与调用它的指令一起决定作用}

一份正确的规则库，如果提示要求“无条件遵循所有条目”，可能比只读其中相关部分更容易出错。反之，相关条目虽然存在，执行提示若没有要求检查其适用条件，也可能从未用上。上下文优化因此至少涉及资产内容、选取方式和使用指令，不能只比较存储里保存了哪些文字。

评价时可以先固定检索结果，观察使用指令改变后的行为，再固定使用指令比较资产内容。这样有助于定位问题，但每种固定都会限制原系统的自由度。若最终目标是完整部署，应再运行允许检索与执行相互作用的整体流程。局部对照帮助解释机制，整体对照判断实际效果，两者的结果应分别报告。

一个条目的存在也不代表它在当前调用中可见。长度截断、排序、模板拼接错误和权限过滤都可能让它消失。若某个候选没有产生预期行为，先检查实际请求是否包含对应规则，再讨论模型是否理解规则。第5章对系统提示的重建审计正是这一原则的实例：存储文件、检索选择和最终可见文本是三层不同证据。

\subsubsection{删除应同时检查依赖和适用范围}

提示条目往往引用共同定义。删除“本章中的任务完成指通过外部验收”这一条，其他规则里“完成后停止”就可能改变含义。表面上只删除了一段重复说明，实际却改变了多个条目的解释前提。因此，维护不只是记录被删文本，还要检查哪些保留内容依赖它。

可以将少量重要定义放在独立的固定契约中，并让可搜索条目引用这些定义；也可以在版本审查时检查交叉引用。无论采用哪一种，都应避免让条目标识与内容分离后失去一致性。条目编号相同而语义已经反转，会使历史有用计数失去解释基础；重大语义变更应留下新版本，而不是继承旧计数后继续累加。

对于被证实有害的条目，还要检查派生内容。它可能已被概括成另一条规则，进入了示例，或被编译到某个模块提示中。只从原位置删除，不保证影响已经消失。撤销应沿来源关系寻找需要复核的产物；若没有保存关系，就如实记录无法完整追踪的范围，不能仅凭存储中搜不到原句便宣布纠正完成。

\subsection{同一任务中的检查点与上下文重置}
\label{sec:ch07-coral}

局部更新之外，CORAL 提供了另一种管理当前上下文的方式\cite{coral-2610}。它让Agent调用工作记忆管理工具，保存任务进度和解题过程中标为已确认的事实；当上下文变长或杂乱时，再调用上下文优化工具，清除会话历史，让下一段从同一系统提示、同一问题和当前工作记忆继续。原文的多个episode属于同一问题的一次连续运行，不能直接当成本书多个独立任务回合。

这种检查点在同题推理片段之间保留信息，本身不证明记忆已经迁移到新任务。上下文恢复与外部动作恢复也须分别定义：重建模型输入不会自动撤销已经发出的工具动作，后者需要环境快照或补偿机制单独承担。这是对接口范围的区分，不是该论文报告的一项外部副作用回滚实验。论文另有监督微调和强化学习来训练检查点策略；训练造成的参数变化，与推理时主动重置上下文是两个层次\cite{coral-2610}。

重置还改变了后续推理能核对什么。作者在附录E指出，移出会话历史也会失去用原始证据复核记忆的渠道，并将为事实携带证据指针列为后续方向，不能把它写成现有出处校验功能\cite{coral-2610}。本书据此建议分开保存完整审计轨迹和下一段实际输入的工作记忆：前者说明信息怎样形成，后者说明Agent当时能看见什么。存下审计记录并不等于模型仍能访问它。

压缩还要保留事实的核验状态。CompAct-Bench先保留原执行者在\cited{3}次筛选运行中均成功的任务，再固定压缩边界的环境状态与剩余动作预算来续跑。作者报告的一条搜索病例中，原轨迹尚未确认的候选，在部分摘要中被写成已找到答案，随后出现过早提交的错误结果\cite{compact-bench-2611}。核对工作记忆时，候选是什么、哪些线索支持或反驳它，以及它是否仍待确认，因而需要一起检查；仅核对实体名称和数值会漏掉这种状态变化。

任务约束则需要另设检查。Salle等人在固定OpenCode框架的实验之外，还用模型评审SWE-chat的真实用户对话；在这批样本中，工作方式与范围限制被判为从后续摘要中丢失的比例高于功能要求\cite{compaction-friendly-2610}。这项对话分析检查约束是否留在摘要中，没有测定丢失对最终任务结果的影响。终局得分、摘要中的约束保留，以及后续行为是否合规，应当分别记录。

\section{搜索中的高分为何需要独立检验}
\label{sec:ch07-selection}

自动优化会反复询问同一批开发任务：“这个候选是不是更好？”当回答带有噪声时，搜索器不仅在选择有效修改，也在选择偶然较高的读数。最终候选的开发集分数已经参与了选择，不能再当作一个未经选择的性能估计。这个问题不依赖优化器是否使用强化学习，人工反复试提示同样会遇到。

用一个可以精确枚举的\textbf{教学例子}来看这种偏差。假定所有候选在每道题上都有相同的二分之一成功概率，各候选、各题结果彼此独立，开发集只有四道题。对每个候选取四题平均分，再挑分数最高者。没有任何候选具有真实优势，但被挑中的最高开发分会随候选数增加。

令单个候选的成功题数为 $X\sim\mathrm{Binomial}(4,1/2)$，比较 $m$ 个候选时，最高开发平均分的期望为
\begin{equation}
\mathbb E\!\left[\frac{\max_i X_i}{4}\right]
=\frac14\sum_{k=0}^{3}\left[1-\Pr(X\leq k)^m\right].
\label{eq:ch07-selection-bias}
\end{equation}
根据二项分布的有限概率逐项求和，可以得到表~\ref{tab:ch07-selection}。随书脚本使用分数精确运算，没有模型调用，也没有随机模拟误差。

\begin{table}[htbp]
\centering
\caption{等能力候选也会出现越来越高的最佳开发分：构造算例}
\label{tab:ch07-selection}
\begin{tabular}{ccc}
\toprule
候选数 & 最高开发分的期望 & 独立新题上的真实成功概率 \\
\midrule
1 & \calcnum{0.5000} & 设定为 $0.5$ \\
2 & \calcnum{0.6367} & 设定为 $0.5$ \\
4 & \calcnum{0.7486} & 设定为 $0.5$ \\
8 & \calcnum{0.8383} & 设定为 $0.5$ \\
\bottomrule
\end{tabular}
\end{table}

这张表不是实际提示优化的性能预测。真实候选往往高度相关，也可能确有能力差异，任务结果未必独立。但例子证明了一个足够重要的事实：即使没有真实改进，选择本身也能产生漂亮的最佳分数。候选搜索次数越多，就越需要将开发期选择与最终评价分开。

如果每轮搜索都查看所谓封存集，再根据其结果生成下一版提示，这批题就已经成为优化输入。即使优化器只能看到一个总分，也可能逐渐适应该总分代表的任务集合。不能通过不公开逐题标签，便宣称封存集从未参与优化。真正需要记录的是哪些反馈被谁在什么时间用于了哪些决策。

\subsection{重复运行有用，但不是唯一的修复办法}

对有随机性的候选进行重复运行，可以减少特定估计量的波动；在相同任务上配对，也能帮助比较候选差异。它们都不能完全消除反复选择同一开发集的偏差。若重复次数很少、候选很多，最终优胜者仍可能受噪声支配；若任务本身没有代表性，重复同一批题更不会创造跨任务覆盖。

可以将预算分成候选生成、初筛、复核和最终验证几个阶段。初筛用较少资源淘汰明显无效者，复核给少数候选更多重复或更广任务，最终只对冻结候选做独立评价。这样做并不要求所有候选耗费相同资源，但资源分配规则应预先说明，选择过程中消耗的全部费用也应计入。

停止规则同样重要。看见一个高分就停，与预算耗尽后选择最高分，是不同搜索策略；某个候选在小样本上超过基线便提前结束，也可能增加偶然成功的机会。若采用顺序决策，应保留每次查看的结果与停止理由，并使用与该规则匹配的统计解释。最终报告不宜只写“搜索十轮”，而省略中间多次重试、手动回退或额外筛选。

\subsection{Pareto 候选保存了什么}

GEPA 不只沿当前平均分最高的提示继续搜索。原文 Algorithm 2 先找出每道开发题上并列最好的候选，取这些集合的并集，再剔除其中被支配的候选，按保留的胜出次数抽取父候选\cite{gepa-2507}。因此，“进入父候选抽样”与“平均分最高”是两个判据。

\begin{table}[htbp]
\centering
\caption{按题胜出与平均分可以选择不同候选：教学构造}
\label{tab:ch07-pareto-example}
\begin{tabular}{lccc}
\toprule
候选 & 开发题一 & 开发题二 & 平均分 \\
\midrule
甲 & \calcnum{1.0} & \calcnum{0.0} & \calcnum{0.5} \\
乙 & \calcnum{0.0} & \calcnum{1.0} & \calcnum{0.5} \\
丙 & \calcnum{0.6} & \calcnum{0.6} & \calcnum{0.6} \\
\bottomrule
\end{tabular}
\end{table}

在表~\ref{tab:ch07-pareto-example} 中，甲、乙各赢一题，丙的均分最高，却没有在任何一题上并列第一。按上述步骤，父候选抽样集合包含甲、乙，不包含丙。丙也没有被甲或乙支配，说明这里的按题胜出集合并不等于完整的多目标非支配集合。丙仍可留在总候选池中，并由 Algorithm 1 的最终平均分判据返回；本例区分的是“从谁继续修改”和“最后交付谁”，不是实际优化性能的预测。

按题最优也可能来自噪声。某个候选偶然答对难题，不代表稳定掌握了相应机制；保留互补方向后，最终产物仍需在未参与选择的任务上评价。这里的比较向量是各题得分。若要讨论费用与质量的 Pareto 曲线，还需把费用纳入目标并实际测量。

组合还会改变模块之间的约定。父候选甲让摘要模块输出结构化实体列表，父候选乙的查询模块却期待自然语言段落；将甲的摘要和乙的查询直接拼起来，可能让两个单独有效的改动同时失效。候选继承记录应包括模块接口与依赖，一次组合也需要重新评价。

\section{离线优化、在线适应与跨模型迁移}
\label{sec:ch07-offline}

最后一个分界是\emph{什么时候优化}。

\textbf{离线（编译期）。}DSPy 的一次 compile 运行：
在训练集与开发集上跑程序、筛演示、选指令，产出一个编译后的程序\cite{dspy-2310}。
GEPA 同属这一侧，它的 rollout 预算花在优化期\cite{gepa-2507}。
离线的好处是预算清晰、可以并行；代价是部署后遇到的新失败不会自动进来。

\textbf{在线（推理期）。}ACE 明确同时支持两侧：
离线优化系统提示、在线适配 Agent 记忆\cite{ace-2510}。
在线这一侧的钱花在每一次任务上，于是``什么时候更新''变成了第 13 章的问题。

\textbf{跨模型}是三者里最容易被说混的一个。
DSPy 的做法是\emph{针对每个基础模型分别编译}并分别报告：
GSM8K 上 GPT-3.5 与 llama2-13b-chat 的数字是两列，不是一列\cite{dspy-2310}。
ACE 原文 §4.5 在其他模型上保持算法和提示模板并重新进行适配，这支持方法可以用于不同模型，却不等于将已生成的同一份上下文直接交给新读取者。原文还将 AppWorld 排行榜中的生产系统作为历史背景参照，并在脚注声明不作为同设置的方法基线\cite{ace-2510}；这也不能证明固定上下文的跨模型迁移。

因此应分别辨认三件事：重新编译或适配、保持算法而换底座、冻结持久物而换读取者。前两种可以展示方法适用范围，第三种才直接检验某份持久物的可移植性。本书 E03c 已完成换模型臂，结论与限制见第5章，回答P3b还需区分模型变化与执行接口变化；不能将该结果预先归给 DSPy 或 ACE。

\subsection{冻结提示迁移的直接证据：GEPA 的另一张表}

不能因为 DSPy 与 ACE 的上述实验没有直接回答固定资产迁移，就忽略其他论文已经提供的相应对照。GEPA 的 Table~2 单列了在 Qwen3-8B 上优化、再不加修改地交给 GPT-4.1 Mini 运行的提示，标作 GEPA-Qwen-Opt。目标模型的基础汇总分为 \cited{53.03}，迁移提示为 \cited{62.03}，两者相差 \calcnum{9.00} 个分数点\cite{gepa-2507}。这是对六个异质基准的论文汇总，不能写成每题成功概率提高九个百分点，更不能将九个分数点改称百分之九的相对增长。

同表直接在目标模型上优化的 GEPA 汇总分为 \cited{65.22}，比迁移版本高 \calcnum{3.19} 个分数点。这个额外参照帮助区分两个问题：旧提示是否优于新模型的起始提示，以及旧提示是否达到重新优化后的表现。迁移版本高于起始基线，并不意味着无需适配也能取得相同收益；反过来，低于直接优化版本，也不意味着这份提示没有可迁移内容。两种部署选择还需分别计算来源搜索与目标适配费用。

这里的“原样”依据作者对提示的报告，不扩张为所有软件和环境条件逐字节相同。本书没有复现该表或核验全部运行轨迹，也不据此补上迁移差值的统计区间。更重要的是，这组对照没有操纵本书 P3b 的两种拟合目标；它为固定提示可以跨模型发挥作用提供了具体案例，却不能替本书完成领域拟合与求解器拟合的比较。相应研究问题需要同协议的拟合目标比较。

\subsection{在线适应的时间顺序就是实验条件}
\label{sec:ch07-online-order}

离线优化通常先完成搜索，再冻结产物进行测试；在线适应则在任务流中边做边更新。后者最容易出现的报告错误，是把使用过当前任务答案的更新应用回当前任务，再将重答结果当作未适应时的首次能力。两种流程都可以研究，但前者测的是未来任务适应，后者包含任务内反馈与修复，不能使用同一名称而不说明顺序。

若目标是测量面向未来任务的在线学习，可以先用当前冻结上下文回答任务 $t$，保存并评分，再释放允许使用的反馈形成下一版本，供任务 $t+1$ 使用。这样当前得分不受事后更新影响。若允许同题重试，就应把首次结果与修复结果分开，并对额外反馈和调用记账。用户真正看到的是哪次答案，也要明确。

任务顺序会改变在线资产。先遇到大量同类问题，可能形成偏向该类的规则；先遇到异常案例，可能让过度保守的限制占据上下文。评价可以预定多个顺序或多条独立任务流，观察结果是否依赖某个方便的顺序。若只运行一个顺序，应将其作为实验条件记录，不用任意打乱后的离线统计假装获得新的独立在线重复。

离线 warmup 也是时间条件。ACE 的 AppWorld 在线结果中，有使用预先适配上下文的版本\cite{ace-2510}。它与从空上下文开始的在线学习并非相同起点。warmup 可以显著影响前几项任务的表现，相关调用与数据也属于系统投入。比较时应分别报告是否 warmup、它使用什么任务、成本是否计入，以及在线阶段是否仍允许继续更新。

\subsubsection{冻结上下文迁移，与重新生成上下文迁移}

当方法换了基础模型仍然有效，至少有两种可能：同一算法能够为新模型重新生成合适资产，或者旧资产本身能够被新模型直接使用。前者证明适配过程具有一定可用性，后者检验持久物的可移植性。再加上模型原有能力差异，最终分数很难仅靠两列总成绩解释。

一个清楚的迁移设计，应为新读取者保留无资产基线，再区分直接使用旧资产与重新适配后的资产。若旧资产需要格式转换，转换规则和费用也要记录；如果转换过程中允许读取目标任务反馈，就已经包含新的学习。不能把重新生成一份内容称作“无成本搬迁同一份经验”。

迁移失败同样需要分解。新模型可能无法执行某种输出格式，可能不支持原工具，也可能只是对提示措辞敏感。先检查接口兼容和资产可见性，再判断内容是否缺乏通用性。否则，一个解析器错误就可能被错误写成“问题域知识无法跨模型”。E03c 当前结果提供有限范围的观察，不能替代针对各类提示程序的迁移测试。

\subsubsection{适应何时应当暂停}

在线系统不是每次失败都需要更新长期上下文。临时网络故障、一次工具超时或不可重复的服务异常，可能不适合形成永久规则；某个失败若来自过时工具说明，则可能需要更新接口知识。应先区分任务机制错误、资料变化与运行故障，再决定修改哪一种持久物。

暂停更新可以由具体可观察事件触发：反馈器版本发生变化、生成的候选无法通过格式检查、关键回归任务显著退化，或资源预算已经耗尽。暂停期间仍可使用最后通过检查的上下文运行，并保存待分析事件。这样做的目的不是保证旧版本永远正确，而是避免在反馈含义已经变化时继续把不确定解释固化成规则。

恢复更新也需要一个边界。应先确认反馈接口和评价口径，再处理积压候选；旧候选若依赖已经失效的工具或资料，应重新审核。一次暂停不应成为之后无预算批量更新的理由。每轮发布都应对应具体候选、具体检查与具体来源，而不是一句“系统已经自行适应”。

\section{对照实验：E04 的三臂与固定查询预算}
\label{sec:ch07-e04}

现在将搜索对象、诊断依据与预算边界合成一个可判的对照。
本书把它交给 E04。先说级别：E04 是\textbf{B 级受限实验}，
截至本章写作时尚未运行；本节写的是对照设计与受限边界，
\textbf{不报告 E04 的实验结果}，也不允许从本节推出“自动搜索优于人工修改”。

\subsection*{三臂}

\begin{itemize}
\sloppy
  \item \texttt{fixed}：系统提示与示例固定，不做任何优化，作为对照基线。
  \item \texttt{human\_edit}：由人根据失败样例改出一版提示，\textbf{只允许改一次}。
  \item \texttt{auto\_search}：按第~\ref{sec:ch07-budget} 节的循环自动生成候选、评估并选择。
\end{itemize}

人工臂和自动臂的可见信息需要对齐。如果改提示的人能看到比自动臂更多的失败样例、更多轮的反馈，
比较就同时改变了编辑方式与可见信息，不能将差异单独归因于人或搜索器。
因此必须预注册三件事：人工臂可见的失败样例集合、修改次数上限、修改者能看到哪些信号。
缺少这些记录会削弱可复核性。当前尚无真实人工编辑产物。若改用模型根据相同材料作一次编辑，应如实命名并冻结新设计，不能据此支持人机优劣结论。

\subsection*{为什么要固定查询预算}

不固定查询预算，\texttt{auto\_search} 可能同时因搜索策略与更多计算而改善；这样的整体方案比较仍有意义，但不能单独归因于搜索方法。本书用的口径是\textbf{固定查询预算}：
三臂在产出最终候选之前，允许向后端发起的评估次数上限相同。
E02 与 E03 提供了按题配对的历史分析，但共同回合上限并未保证完整搜索与来源成本相等，不能说已经验证了这里的等预算设计。

固定预算正是比较样本效率的一种方式：在相同可用查询数下，观察谁取得更好的留出表现。还可以预定多个预算点，画出质量随查询数变化的曲线。单一预算点不能覆盖所有资源区间，但不应说固定查询预算“压住了”样本效率优势。查询数之外仍须记录反思、候选生成、验证和服务调用的实际费用。

\subsection*{固定候选的外部评测}

自动搜索最大的风险不是搜得差，而是搜到的东西只在本轮优化集上好看。
E04 的处置是把三臂产出的最终候选，放到\textbf{同一批未参与优化的封存题}上评，
并把优化集收益与封存集收益分开报。优化集改善而封存集点估计接近零，是需要检查泛化的信号；若要认定收益小于实际重要阈值，需预定阈值并有足够精度，不能将“不显著”当作相等。

\subsection*{受限边界}

\begin{itemize}
\sloppy
  \item \textbf{在本节等长度预算问题下，三臂使用相同上限。}否则结果同时包含可用上下文资源的差异，
        无法单独解释为选择策略的作用。第~\ref{sec:ch07-delta} 节那次塌缩说明
        长度本身就是一个会自己失控的量，它不能被当成免费资源。
  \item \textbf{成本要分别记录任务运行与模型资源。}rollout 数与 token 数不能互换，
        一次带反思的 rollout 与一次纯生成的 rollout 单次成本不保证相等
        （见第~\ref{sec:ch07-budget} 节的口径更正）。只报前者会把结论锁在样本效率这一维上。
  \item \textbf{反馈层级必须写进臂的定义。}E02 在本书任务集上测到四种反馈来源里
        位置最低的是内在自评（\measured{0.6125}），五个非基线臂全部低于无记忆基线
        \measured{0.8125}（运行 \path{E02-20260924T153511Z-421eac}）。
        E04若采用自评式反馈，应另外设置外部反馈对照；E02读数不是新实验的性能上限。
  \item \textbf{候选差异必须留档。}只报最终候选等于只报搜索结果，
        读者无法判断“这一步改动对应哪一个诊断”。交付物里要有候选差异与搜索账本。
\end{itemize}

% ---------- 5. 条件变化 ----------

\subsection{条件变化与结果记录}
\label{sec:ch07-conditions}
任务族、反馈器、读取模型与提示长度上限，是需要分别登记的条件。第5章的族留出与冻结资产换读取者，回答的是不同迁移问题；E02的反馈来源比较也不能替代对同一反馈器噪声的测量。新设计应检查共享记忆、任务族与序列依赖，再确定重采样单位，不能预设所有实验都按题自助。

提示优化还可能把一次偶然判断写进长期规则。检验这条路径，需要关联反馈事件、候选差异和后续行为；开头的塌缩案例只展示了内容丢失，并未单独验证噪声持久化。重复运行或稳健选择可以减少部分波动，最终仍须用未参与选择的任务评价。

\label{sec:ch07-result-card}
E04的结果卡应能恢复三件事：系统看过哪些失败，允许改变哪些行为，以及为什么选择最后一版。冻结卡先登记模型、工具、数据划分、起始提示、可改模块、预算与长度上限；搜索记录再保存父版本、诊断、差异、开发期读数与累计成本，包括被拒绝的候选。最终选定版本应与封存运行实际加载的版本一致。

优化集、封存集和预定分组分别报告，附任务数、重复数、缺失与不确定性单位。若区间很宽，泛化仍未确定；若区间足够窄且落在预定的实际无意义范围内，才有依据说收益很小。失败分析则检查候选是否重复答案模板、增加无效调用、删掉例外，或只对特定题型有效。评价来自实际请求、响应和状态变化。

\section{选择原则}
\label{sec:ch07-choice}

\begin{enumerate}
  \item \textbf{先问删掉的部分谁负责。}整体重写应保留前后版本与差异。选增量 delta，或给重写加一条``本次删除了什么''的记录。
  \item \textbf{把诊断与策展分开调用。}Reflector 提炼经验，Curator 提出结构化更新，合并程序应用更新。
        分离角色便于定位各阶段；合并调用也可要求结构化记录，但要实际检查记录是否完整。
  \item \textbf{报告预算口径，不要只报 rollout 数。}GEPA 的 \cited{1/35}
        涉及最佳候选出现时点，不能代替完整优化的预算比；还须分别报告延迟、token 与实际费用。
  \item \textbf{把有用/有害计数当机制量，不要当效果量。}
        它说明条目被标记为有用或有害的次数，不说明系统因此变好多少。
  \item \textbf{跨模型主张必须说明是重新编译还是直接换读取者。}
        重新适配与冻结迁移是两种评价操作，不与第 5 章的领域拟合、求解器拟合一一对应；任一种拟合目标都可以接受这两种评价。
\end{enumerate}

% ---------- 7. 反例 ----------
\begin{pitfallbox}
\textbf{反例一：把提示变短当成收敛。}
ACE 讨论了短而通用的提示可能遗漏细节的风险\cite{ace-2510}。短提示也可能足够有效，长度稳定与质量收敛都需检查，不能由变短单独诊断搜索失败。

\textbf{反例二：把上下文变长当成信息在增长。}
AppWorld 上第 60 步的 \cited{18{,}282} token 到第 61 步变成 \cited{122} token，
准确率从 \cited{66.7} 掉到 \cited{57.1}，低于不做适配的 \cited{63.7}\cite{ace-2510}。
长度是双向的：塌缩之前的那份长度里，有多少是真正被用到的，
在那一步之前无法从长度本身判断。

\textbf{反例三：把优化集上的提升当成能力变化。}
优化集收益与封存集收益不一致，需要检查选择偏差、任务分布与估计精度；这一现象本身不能唯一归因于记住题目。
表~\ref{tab:ch07-ace} 里的数字按原论文各自条件报告，
本书未复现、也未在按族留出的划分上看过它们——
\emph{在看到族留出数字之前，这些表格支持的是所报告条件下的表现，不独自证明任意跨族迁移}。

\textbf{反例四：把``少 rollout''当成``更便宜''。}
见 §\ref{sec:ch07-budget} 的原则框：约 \cited{1/35} 涉及事后找到的最佳候选出现时点，
不是完整搜索费用比，单次 rollout 也不是统一成本单位。
\end{pitfallbox}

\section*{习题}
\begin{enumerate}[label=7-\arabic*.,ref=7-\arabic*]
  \item\label{ex:ch07-1} 选一个你正在用的提示优化流程，写出它的一次更新动作：
        是整体重写还是增量 delta？如果是前者，请估计一次更新里平均有多少条旧内容消失，
        并说明你打算怎么测这个数。
  \item\label{ex:ch07-2} 表~\ref{tab:ch07-cost} 里离线场景的延迟从 \cited{53898} 秒降到 \cited{9517} 秒。
        若要区分更新方式与候选评估调用量的作用，你会怎样构造交叉对照，哪些条件必须固定？
        若两因素存在交互，为什么不能把这 \cited{82.3\%} 唯一分成两份贡献？
  \item\label{ex:ch07-3} 本章说``有用/有害计数是机制量不是效果量''。
        请设计带对照的试验估计条目对结果的作用，说明需要额外记录什么，以及为何计数本身仍不是效果量。
  \item\label{ex:ch07-4} ACE 在多个基础模型上运行相同适配方法。
        要判断同一份已生成上下文是否可移植，你会如何冻结资产并设置新读取者基线？
\end{enumerate}

\section*{常见误解}
\begin{sloppypar}
\begin{itemize}
\sloppy
  \item \textbf{``上下文适配就是提示工程换了个名字''。}
        差别在搜索对象。DSPy 优化模块指令、示例等可编译内容并留下编译产物\cite{dspy-2310}，
        ACE 优化的是一组带计数的条目集合\cite{ace-2510}。
        手写提示同样会增删，也应保存版本与评价依据。
  \item \textbf{``反思器看一下就能写出好洞见''。}
        E02 在本书任务集上测到自评配置的封存题回合通过率 \measured{0.6125} 低于环境结果的
        \measured{0.7000}。该差异区间跨零，且评审请求并不继承求解历史，不能由点估计证明角色因果。
  \item \textbf{``长上下文一定贵''。}
        ACE 报了 prompt caching 下 \cited{91.8\%} 的输入 token 来自缓存，
        输入 token 费用相对按原始长度计费降 \cited{82.6\%}\cite{ace-2510}。
        原始 token 数、缓存 token 数与实际费用不是同一个量，成本模型要先分清这几项。
  \item \textbf{``换模型能直接用，说明上下文是通用的''。}
        也可能是任务本身没区分度。判别方法是补一个换模型的无记忆基线——
        这正是 E03c 里 \texttt{base\_cross} 这一臂存在的原因。
\end{itemize}
\end{sloppypar}

\section*{进一步阅读}
\begin{itemize}
\sloppy
  \item \cite{ace-2510}：context collapse 与 brevity bias 的原始案例，
        Figure~2 是本章开头那三个数字的出处，Table~4 是本章成本表的出处。
        读 §2.2 时注意它把塌缩描述为整体重写的普遍风险，而非某个方法的缺陷。
  \item \cite{gepa-2507}：reflective prompt evolution 与 GRPO 的对照，
        摘要增益需结合 Table~1/2 的差值口径；IFBench 最佳候选出现时点与总优化预算须分开读取。
        读它时先对齐预算口径；Table~2 另有冻结提示直接迁移的对照，不能与重新适配合并解释。
  \item \cite{dspy-2310}：signature / module / teleprompter 三抽象与编译三阶段，
        Table~1、Table~2 是各模型分别报告的样本——``换模型就重新编译''这一说的出处。
  \item 本书第 4 章：历史反馈配置的边界与 E02 六臂实测，
        自评臂 \measured{0.6125} 与环境臂 \measured{0.7000} 的对照在那里有完整区间。
  \item 本书第 5 章：两种拟合目标与 E03c 的换模型结果及未决问题，
        本章 §\ref{sec:ch07-offline} 的迁移问题在那里有已完成运行与明确未决范围。
\end{itemize}
